\section{Twierdzenie o niedefiniowalności prawdy (Tarski)}\label{sec:tarski}
\begin{twierdzenie}[Tarski: wersja semantyczna i syntaktyczna]\label{res:10:1}
Nie istnieje formuła arytmetyczna $\operatorname{Tr}(x)$ definiująca prawdę wszystkich zdań arytmetycznych w $\mathcal N$, czyli spełniająca dla każdego zdania $\varphi$
\[
\mathcal N\models\operatorname{Tr}(\ulcorner\varphi\urcorner)\quad\Longleftrightarrow\quad\mathcal N\models\varphi.
\]
Ponadto żadna niesprzeczna teoria $T\supseteq Q$ nie dowodzi wszystkich równoważności $\varphi\leftrightarrow\operatorname{Tr}(\ulcorner\varphi\urcorner)$ dla jednej takiej formuły $\operatorname{Tr}$. Efektywność $T$ nie jest potrzebna.
\end{twierdzenie}
\begin{proof}
Lemat~\ref{res:8:1} zastosowany do $\psi(x)=\neg\operatorname{Tr}(x)$ daje zdanie $L$ z $Q\vdash L\leftrightarrow\neg\operatorname{Tr}(\ulcorner L\urcorner)$. Ponieważ $\mathcal N\models Q$, równoważność jest prawdziwa w $\mathcal N$. Gdyby $\operatorname{Tr}$ definiowała prawdę, dla tego samego $L$ mielibyśmy także $\mathcal N\models L\leftrightarrow\operatorname{Tr}(\ulcorner L\urcorner)$. Daje to $\mathcal N\models L\leftrightarrow\neg L$, co jest niemożliwe.

W wersji syntaktycznej obie równoważności byłyby twierdzeniami $T$. W klasycznej logice z $L\leftrightarrow\neg L$ otrzymujemy $\neg L$ (zakładając $L$, otrzymujemy sprzeczność), a potem $L$ z drugiej implikacji. Teoria byłaby sprzeczna.
\end{proof}

\begin{uwaga}\label{res:10:2}
Porównanie z Tw.~\ref{res:9:2}: dowodliwość $\operatorname{Bew}_T$ jest definiowalna (rekurencyjnie przeliczalna), lecz prawda nie jest. Zdanie Gödla wykorzystuje różnicę między „dowodliwe” a „prawdziwe”.
\end{uwaga}

% Rozdziały 11--13, strony 20--24.
